\section{Global Production Systems and Climatic Conditions}

Honey quality standards, including the Apimondia position, were
developed on the basis of the European temperate climate and
frame beekeeping, yet more than 70\% of global honey production
comes from regions that are climatically, technologically, and
culturally different from
Europe~\cite{garcia2025,eudir2001110}. Applying uniform morphological and microbiological criteria to honeys from different production systems is methodologically limited and may create disproportionate assessment outcomes for producers operating under different climatic and technological conditions.~\cite{phiri2022,tadele2025}.

%,----------------------------------------------------------
\subsection{Global Landscape of Apicultural Production}
%,----------------------------------------------------------

According to FAO data from 2024, the global number of managed
colonies exceeded 101.7~million, with Asia dominating
at a 44.4\% share, and Europe holding 24.7\%~\cite{destatis2026}.
Honey production in 2022 reached 1.83~million tonnes, with
Asia accounting for 48.2\%, Europe for 22.8\%, the Americas
for 18.5\%, and Africa for 8.5\%~\cite{phiri2022}. Adopting
the capping criterion as the sole determinant of maturity in
practice favors European honeys and discriminates against honeys
from tropical and traditional
systems~\cite{tadele2025,phiri2022}.


% \subsection{African production: climatic humidity as a technological determinant}


Relative humidity in the tropical zone of Africa is typically
70--90\% RH, exceeding 95\% RH during the rainy season.
Under such conditions, bees are not biologically capable of
reducing honey moisture to $\leq 20\%$ solely through natural
hive ventilation~\cite{tadele2025}. The beekeeper faces a choice:
early harvesting and technological dehumidification, or the risk
of in situ fermentation or plundering by pests. Tadele et al.~\cite{tadele2025}
classify dehumidification as a \emph{technological necessity},
not an intent to adulterate, emphasizing that the alternative
is crop loss or fermented product. Studies show that honey
subjected to proper dehumidification retains full enzymatic
and antioxidant profile~\cite{lastriyanto2020,spel3upm}.

\subsubsection{African Production: Climatic Humidity and Traditional Hive Systems}

In Ethiopia, 95.5\% of hives are traditional, with only 0.2\%
being frame hives~\cite{ethaau2007}. Paradoxically, the study
by Gobessa et al.~\cite{sciendo2012} demonstrated that honey
from traditional hives is characterized by \emph{significantly
lower} moisture than from modern hives ($p < 0.001$), as
harvesting occurs once or twice per year, giving more time
for natural evaporation. At the same time, traditional hives
make frame inspection and assessment of the degree of capping
impossible without destroying the combs; pplying a frame-based morphological criterion to these systems illustrates the limited transferability of maturity indicators derived mainly from temperate-climate European production systems\cite{tadele2025}.

%,----------------------------------------------------------
\subsection{Honeys of Stingless Tropical Bees}
%,----------------------------------------------------------

Stingless bees (Meliponini, $>500$~species) produce honey with
moisture typically of 25--35\% and $a_w = 0.60$--$0.86$,
a different sugar profile, and higher acidity than
\textit{Apis mellifera} honeys~\cite{stinglessyeastbrazil2021,stinglesswjava2024}.
Meliponini do not possess nectar dehydration mechanisms
comparable to the European bee: smaller honey crops and less
developed hive ventilation are biological characteristics,
not production deficiencies. A study of 17 species
from Brazil~\cite{stinglessyeastbrazil2021} demonstrated a range
of $a_w = 0.60$--$0.86$, with osmophilic yeasts isolated
from these honeys characterized by high osmotolerance and low
fermentative capacity -- biological adaptation to an environment
with elevated $a_w$. Meliponini honeys require separate
standardization; assessing their quality on the basis of
\textit{Apis mellifera} standards is biologically
unjustified~\cite{stinglessyeastbrazil2021,stinglesswjava2024}.

%,----------------------------------------------------------
\subsection{Specific Microclimates and Weather Conditions in the Temperate Zone}

Honey extraction with moisture levels that pose a real risk of
fermentation also occurs in the temperate zone, including in
Europe. Many areas from which honey is harvested are characterized
by a specific microclimate with persistently elevated air humidity.
Weather factors such as rainfall lasting for periods of several
days to several weeks influence elevated water content in combs.
This is precisely why the Apimondia position states that a justified
beekeeping practice is to leave honey after extraction in open
vessels for the purpose of evaporating excess water. Extraction of
honey with elevated water content is therefore a fact, as is human
intervention aimed at removing excess water for the purpose of
microbiological stabilization. Commonly practiced by European
beekeepers is not only periodic mixing of honey to increase the rate
of evaporation, but also the use of various technical devices for
this purpose (e.g., rotating discs that increase the evaporation
surface). Technical development and process automation are visible
at all stages of honey extraction and serve not only to increase
efficiency but also effectively enhance product safety.
Unquestionably, however, they cannot be instruments of honey
adulteration. Therefore, clear, accessible, and verified criteria
contained in normative acts should underpin the assessment of such
activities, and within justified scope, prohibited practices should
be precisely specified (e.g., the use of nanofiltration, ion exchange
resins, vacuum dehydration, etc.).
%,----------------------------------------------------------
\subsection{Comparison of Production Systems}
%,----------------------------------------------------------

Table~\ref{tab:hive_systems} summarizes the key characteristics of
the three globally dominant hive types. Only the Langstroth hive
provides full control over honey ripening; KTBH and traditional
hive systems are structurally limited in this respect, meaning that
honeys from these systems may fail to meet the capping criterion
while complying with all good beekeeping
practices~\cite{tadele2025,sciendo2012,lrrd2013}.

\begin{table}[H]
\centering
\caption{Comparison of beekeeping production systems with respect to
         honey maturity control. Source: own compilation
         based on~\cite{tadele2025,sciendo2012,lrrd2013}.}
\label{tab:hive_systems}
\begin{adjustwidth}{-\extralength}{0cm}
\small
\begin{tabularx}{\fulllength}{p{3.5cm}p{3.5cm}p{3.5cm}p{3.5cm}}
\hline
\textbf{Feature} &
\textbf{Langstroth hive} &
\textbf{Kenya Top Bar Hive} &
\textbf{Traditional hive} \\
\hline
Capping control
    & Full (frame inspection)
    & Partial
    & None (destructive) \\
\hline
In situ dehumidification
    & Yes
    & Limited
    & None \\
\hline
Average yield [kg/hive/year]
    & 20--40
    & 15--26~\cite{lrrd2013}
    & 5--15~\cite{tadele2025} \\
\hline
Typical honey moisture
    & 16--19\%
    & 17--20\%
    & 17--23\% \\
\hline
Risk of yeast colonization
    & Low
    & Moderate
    & High \\
\hline
Global use
    & Americas, Australia, Europe
    & East Africa
    & Africa, S.E. Asia \\
\hline
\end{tabularx}
\end{adjustwidth}
\end{table}

%,----------------------------------------------------------
\subsection{Limitations of the Morphological Capping Criterion}
%,----------------------------------------------------------

Based on the above analysis, the criterion of full capping as the
sole determinant of honey maturity is \emph{eurocentric} and
\emph{technically inadequate} for five reasons: (1)~the climatic
impossibility of fully dehydrating nectar under conditions of
70--90\%~RH humidity~\cite{tadele2025}; (2)~the technological
inaccessibility of capping assessment in traditional hives without
comb destruction~\cite{sciendo2012}; (3)~the biological distinctness
of Meliponini honeys that precludes direct application of
\textit{Apis} standards~\cite{stinglessyeastbrazil2021};
(4)~the empirical inconsistency between the degree of capping and
the microbiological stability of honey~\cite{zhang2021,guo2020};
(5)~the impossibility of gathering, archiving, and verifying
evidence, and of post factum proof from the degree of capping.
It is necessary to adopt an integrated model for assessing
functional maturity based on measurable parameters ($a_w$),
not on a morphological
criterion~\cite{tadele2025,codex_stan12,ruegg1981}.
